\section{\textbf{Discussion, Limitations, and Future Work}}

AGI-Former is intended to convert the general objective of integrated multimodal agency into an explicit and testable neural architecture. Its central departure from externally orchestrated multimodal systems is architectural: modality-specific encoders, causal thought--action representations, cross-modal fusion, and output generation are trained as components of one end-to-end system. This does not imply that every component must share the same parameters or operate at the same temporal resolution. Rather, the proposal requires the components to participate in a common learned decision process whose internal representations and outputs can be optimized jointly.

This distinction matters because a collection of capable specialist models does not automatically constitute an integrated agent. Systems such as HuggingGPT demonstrate that a language model can plan tasks and route requests among external models \cite{shen2023hugginggpt}. Such orchestration is useful, but its interfaces, routing rules, representation conversions, and failure-handling logic are partly imposed outside the learned model. AGI-Former instead asks whether multimodal perception and action can be coupled through a persistent causal state and learned end-to-end. The resulting architecture should therefore be judged by the controlled comparisons in Section~VII, not by the presence of multiple modalities alone.

\subsection{\textbf{Interpretation of the Architectural Variants}}

The disjoint, joint, and hierarchical variants address different output structures. The disjoint architecture is appropriate when modalities, tools, effectors, or regions require distinct vocabularies and independently constrained outputs. Its specialization may improve modularity, diagnosis, and failure containment, but separately generated streams may conflict. The joint architecture produces a unified thought--action sequence and can represent dependencies among outputs directly. Its shared sequence may improve global consistency, although it can also create a bottleneck when several actions must occur concurrently or at substantially different rates.

The hierarchical architecture combines these properties. A joint global model can express a common objective, allocation, or plan, while disjoint regional models translate that supervisory state into locally valid actions. The two output forms are therefore complementary rather than successive claims that one is universally superior. Their relative value depends on action semantics, concurrency, communication constraints, and the scale of the environment. This dependence is why Section~VII treats output organization as an empirical hypothesis.

The same caution applies to Unified Attention. The UTR results motivate Unified MHA and Unified Cross-MHA as candidate operations for integrating self- and cross-attentive computation \cite{zohourian2026utr}; however, transferring those operations from image classification and captioning to multimodal agency is an architectural hypothesis. Only matched ablations can determine whether their benefit persists in the proposed setting.

\subsection{\textbf{Primary Limitation: Integrated Training Data}}

The most immediate limitation is the lack of a coherent training dataset for the complete AGI-Former objective. Existing multimodal datasets usually supervise bounded relationships: an image and caption, a video and description, an audio--visual event, an instruction and robot trajectory, or a question and answer. Multimodal learning research has long identified representation, alignment, translation, fusion, and co-learning as distinct technical challenges \cite{baltrusaitis2019multimodal}. AGI-Former requires these relationships to coexist across time within episodes containing synchronized sensory streams, latent or observable state transitions, thought or planning traces where available, executable actions, environmental consequences, uncertainty, and corrective feedback.

No claim is made here that currently available datasets jointly provide all of these signals at sufficient diversity and scale. Embodied datasets and environments associated with VIMA, PaLM-E, and RT-2 provide useful starting points for multimodal prompting, embodied representation, and vision--language--action learning \cite{jiang2023vima,driess2023palme,zitkovich2023rt2}. Nevertheless, they do not by themselves resolve the broader data requirements of an architecture spanning language, images, audio, video, persistent reasoning, heterogeneous actions, and hierarchical multi-agent feedback.

One possible mitigation is to use videos with rich scientific, technical, or procedural content as surrogate multimodal episodes. A high-quality scientific video may jointly contain temporally ordered visual evidence, speech, text, demonstrations, causal explanations, and observable state changes. Transcripts, diagrams, temporal segments, and demonstrated procedures could be transformed into aligned training targets. Synthetic questions, predictions, action descriptions, and counterfactuals could further enrich an episode.

However, video is not an automatic replacement for interaction data. Passive recordings often omit the decision alternatives considered by an actor, the unobserved state of the environment, the exact action interface, negative outcomes, uncertainty, and feedback following counterfactual choices. Automatically generated annotations may also reproduce errors from the models that generate them. Scientific video should therefore be treated as a promising source for pretraining and dataset construction, not as an established solution. Determining which video domains, annotation procedures, and quality controls produce transferable agency is left to future work.

\subsection{\textbf{Training Cost and Scalability}}

The second major limitation is computational cost. Multiple high-capacity encoders, persistent causal decoding, repeated cross-modal attention, and long action histories may make full end-to-end training substantially more expensive than training or adapting a single-modality model. Video and audio increase sequence length, while the hierarchical architecture adds communication and coordination overhead. Parameter count alone therefore understates the relevant cost; memory, training FLOPs, data preprocessing, communication, inference latency, and energy must also be measured.

Future work will investigate cost reduction rather than assuming that the largest implementation is necessary. Candidate strategies include initializing encoders from pretrained modality models, freezing and progressively unfreezing components, parameter-efficient adaptation, distillation, temporal and spatial token compression, sparse or gated modality activation, cached representations, curriculum learning, and staged training of regional and global models. Perceiver IO and Gato illustrate complementary approaches to structured multimodal input--output processing and generalist sequence modeling \cite{jaegle2022perceiver,reed2022generalist}; AGI-Former experiments should determine which efficiencies can be adopted without removing the proposed action-conditioned integration.

Initial experiments should therefore be deliberately bounded. A small number of modalities, restricted token budgets, compact encoders, and controlled simulated environments can test the architectural mechanisms before expensive scaling. Positive evidence at small scale would not establish that the architecture scales, but negative or neutral ablations could prevent unnecessary large-scale training. Subsequent experiments can increase model capacity, episode length, modality diversity, and the number of regional agents while reporting the matched budgets and Pareto comparisons defined in Section~VII.

\subsection{\textbf{Additional Limitations}}

Several further limitations remain. First, modalities may differ greatly in information density, noise, sampling rate, and dataset size; without careful normalization and sampling, a dominant modality may suppress weaker streams. Second, joint optimization can produce catastrophic interference, in which improvement on one task or modality degrades another. Third, persistent thought--action histories can accumulate errors over long horizons and may consume increasing context without reliable state compression. Fourth, generated action tokens are useful only when grounded in valid action spaces and verified against physical, software, legal, or safety constraints.

Fifth, hierarchical control introduces new failure modes: inaccurate global plans may propagate to several regions, regional messages may be stale or deceptive, and communication loss may divide the system's state. The fallback policies described in Section~VI reduce but do not eliminate these risks. Finally, interpretability remains limited. Attention weights or readable language tokens should not be treated as complete explanations of a decision; causal interventions, execution traces, and controlled counterfactuals are needed to evaluate what information affected an action.

These limitations also define the future-work program. The next steps are to construct coherent multimodal episodes from existing embodied resources and carefully selected scientific videos; define aligned sensory, state, thought, action, and feedback schemas; implement compact disjoint and joint prototypes; perform the mechanism-level ablations in Section~VII; and scale only those components supported by controlled evidence. Later work can examine hierarchical multi-agent deployment, continual adaptation, missing-modality robustness, uncertainty-aware supervision, and safe human intervention.

AGI-Former consequently remains a proposed architecture and research program, not a demonstration of AGI. Even successful experiments would establish evidence for integrated multimodal agency under specified conditions, rather than universal intelligence, consciousness, or superintelligence. This boundary is consistent with the distinction between consciousness and AGI and with the characterization of AGI as integrated multimodal agency developed in prior work \cite{zohourian2026agi}. The contribution of the present paper is to make that characterization implementable, comparable, and falsifiable.
